\chapter{كيف تتعلم الآلة}
% full version: chapters/06_dofa.tex

في كل الدوائر حتى الآن كان المهندس هو من ينتقي قوة الوصلات. وفي الدماغ الحي لا مهندس. يجب أن تُضبط الوصلات بنفسها، في أثناء العمل، وعلى نحو يجعل الآلة تفعل شيئًا مفيدًا. هذا هو التعلم.

القيد الأساسي صارم جدًا. كل وصلة تغيّر نفسها بنفسها، ولا يمكنها أن تعرف إلا ما يجري قربها: هل يعمل المرسِل، وهل يعمل المستقبِل، وأي المواد تصل إليها. لا أحد يستطيع أن يرسل إلى الوصلة تصحيحًا خاصًا بها. وهذا يستبعد الطريقة الرئيسية لتعليم الشبكات العصبية الاصطناعية، حيث يجري الخطأ في الشبكة بالاتجاه المعاكس وتنال كل وصلة تصحيحها الخاص.

\section*{معًا يعني مترابطين}

أبسط قاعدة: إذا عمل المرسِل والمستقبِل معًا، قويت الوصلة بينهما. إنها محلية ولا تحتاج إلى ساعة. وقاعدة كهذه، إذا أضيف إليها قليل من النسيان، تجد في تدفق الإشارات أهم ما فيه: ما يتغير أكثر من غيره. وثمة صيغة تنظر إلى توقيت النقرات: تقوى الوصلة إذا نقر المرسِل \emph{قبل} المستقبِل، أي كان يمكن أن يكون السبب، وتضعف إذا نقر بعده.

لكن لهذه القواعد حدًا واحدًا: إنها تعلّم ما هو \emph{متكرر}، لا ما هو \emph{مفيد}. الوصلة التي لا ترى إلا جيرانها لا تعرف هل جلب الفعل عصيرًا أم ألمًا.

\section*{الإشارة الثالثة}

يأتي بها الدوبامين، محطة ”أفضل مما كان متوقعًا“. لنجعل الوصلة لا تتغير إلا إذا اجتمعت ثلاثة أمور: عمل المرسِل، وعمل المستقبِل، ووصل الدوبامين.

\pencilfig{6}{\pencilart{6}}{تقوى الوصلة حين يعمل العصبونان معًا وتصل الإشارة الثالثة: ”جاءت النتيجة أفضل مما انتظرنا“.}

لكن هنا صعوبة: المكافأة لا تأتي فورًا، بل بعد ثانية أو ثانيتين. وحينها يكون المرسِل والمستقبِل قد صمتا منذ زمن. تحتاج الوصلة إلى ذاكرة بأنها شاركت. وهذه الذاكرة موجودة: العمل المشترك يترك في الوصلة علامة، كملصق ينفصل وحده خلال ثانية أو اثنتين. إذا وصل الدوبامين والملصق لا يزال عالقًا، تغيرت الوصلة. وإذا لم يصل، زالت العلامة بلا أثر. وهذا يحل مسألة العنونة: الدوبامين واحد للجميع، لكن لا يتغير إلا الوصلات التي عُلّق عليها الملصق.

\pencilfig{106}{%
\foreach \y/\d/\t in {0/0.9/{في الوقت: الوصلة تقوى},-1.3/3.3/{متأخر: لا شيء}}{
  \draw[pen] (0,\y) -- (4.6,\y);
  \foreach \s in {0.25,0.33}{\draw[pcBlue, line width=0.8pt] (\s,\y) -- (\s,\y+0.3);}
  \fill[pcAmber, opacity=0.25] (0.35,\y) -- plot[domain=0.35:4.5, samples=30] (\x,{\y+0.8*exp(-(\x-0.35)/0.8)}) -- (4.5,\y) -- cycle;
  \draw[penlite=pcAmber] plot[domain=0.35:4.5, samples=30] (\x,{\y+0.8*exp(-(\x-0.35)/0.8)});
  \draw[pcAmber!80!black, line width=0.9pt] (\d-0.1,\y+0.95) -- (\d+0.07,\y+0.62) -- (\d-0.07,\y+0.6) -- (\d+0.1,\y+0.22);
  \draw[arr=pcAmber!80!black] (\d+0.09,\y+0.27) -- (\d+0.11,\y+0.12);
  \node[lbl, anchor=west] at (4.7,\y+0.3) {\t};}
\node[lbl, text=pcBlue, anchor=south west] at (0.1,0.85) {عملا معًا};
\node[lbl, text=pcAmber!80!black, anchor=south] at (2.3,0.35) {العلامة تذوب};
\foreach \x/\l in {0.35/0,1.95/1,3.55/2}{\draw[pcGray, line width=0.5pt] (\x,-1.3) -- (\x,-1.38); \node[lbl, anchor=north] at (\x,-1.38) {\foreignlanguage{english}{\l~s}};}
}{العمل المشترك يترك في الوصلة علامة تذوب خلال ثانية أو اثنتين. الدوبامين لا يغيّر إلا الوصلات التي تحمل العلامة.}

\section*{الضجيج بحثًا}

لماذا تقود هذه القاعدة إلى الأفضل لا إلى أي مكان؟ تخيّل أنك تبحث عن قمة تل وعيناك معصوبتان. تخطو خطوة عشوائية، فيصيح أحدهم ”أعلى!“ أو ”أدنى!“. إذا كانت ”أعلى“، ثبّتّ الخطوة. خطوة بعد خطوة تصعد، دون أن تنظر إلى الخريطة مرة. والدماغ يفعل الشيء نفسه: العصبونات ترتعش قليلًا ارتعاشًا عشوائيًا، والدوبامين يخبر هل صار الوضع أفضل، والوصلات التي شاركت تنزاح في الاتجاه المفيد. الضجيج الذي كان يعيق نقل الرسائل يصبح هنا بحثًا.

والآن يتضح لماذا لا يخبر الدوبامين عن المكافأة، بل عن \emph{المفاجأة}. لو كان يصيح ”عصير!“ بعد كل نجاح، لثبّت معه كل ما كانت الشبكة تفعله، الصحيح والخاطئ. وفي نموذجنا تضخّم شبكة كهذه وصلاتها فعلًا آلاف المرات، وأحيانًا تعلق إلى الأبد عند أسوأ خيار. طرح التوقع يجعل التعلم نزيهًا.

\section*{ثمن السلك الواحد}

للتعلم بالبث ثمن. إشارة واحدة للجميع، وفيها يختلط ضجيج كل العصبونات التي تؤثر في النتيجة. كلما كبرت الشبكة، طال الوقت الذي تحتاجه كل وصلة لتستخلص حصتها. لذلك يبدو أن الدوبامين لا يعلّم إلا دوائر صغيرة نسبيًا تختار الأفعال، أما القشرة الهائلة فتتعلم أساسًا بقواعد محلية.

\section*{من يحسب المفاجأة}

يبقى سؤال: كيف تحسب عصبونات الدوبامين ”أفضل مما كان متوقعًا“؟ نحتاج إلى ناقد: مجموعة عصبونات تقدّر طوال الوقت كم من الخير لا يزال قادمًا. المفاجأة هي المكافأة المتلقاة زائد مقدار قفزة التوقع. أضاء المصباح، فقفز التوقع: رشقة. وصل العصير الموعود: المكافأة موجودة، لكن التوقع استُنفد في اللحظة نفسها، فلا مفاجأة. لم يصل العصير: انهار التوقع بلا مقابل، ففجوة. الحالات الثلاث من الفصل الأول تخرج من تلقاء نفسها. أن الدوبامين يتصرف هكذا بالضبط أمر مقيس. أما كيف يحسب الدماغ هذه المفاجأة بالضبط فلا يزال فرضية.

\fullversion{6}
